\chapter{Creare un Plugin WordPress}

\section{Introduzione}

Un \textbf{plugin} di WordPress è un'estensione che aggiunge funzionalità al CMS. In questo capitolo vedremo come strutturare e implementare un semplice plugin, esaminando sia la struttura dei file sia i principali meccanismi utilizzati da WordPress per interfacciarsi con i plugin (azioni, filtri, hook).

\section{Struttura di file e directory}

Un plugin WordPress è tipicamente organizzato in questo modo:

\begin{itemize}
    \item \textbf{Cartella del plugin}: es.\ \texttt{myplugin}.
    \item \textbf{File principale del plugin}: es.\ \texttt{myplugin.php}.
    \item \textbf{Cartella plugin di WordPress}: \texttt{wp-content/plugins}.
\end{itemize}

Il percorso completo del file principale sarà:

\begin{verbatim}
wp-content/plugins/myplugin/myplugin.php
\end{verbatim}

\section{Header del file principale}

Il file principale deve contenere un header con le informazioni del plugin. WordPress utilizza questi metadati per riconoscere il plugin e mostrarne le informazioni nella bacheca:

\begin{lstlisting}[language=PHP]
<?php
/*
Plugin Name: My Plugin
Version: 1.0
Description: Descrizione del plugin
Author: Nome Cognome
Author URI: http://sito.com
*/
\end{lstlisting}

I campi principali dell'header sono:

\begin{itemize}
    \item \texttt{Plugin Name}: il nome del plugin.
    \item \texttt{Version}: la versione del plugin.
    \item \texttt{Description}: una breve descrizione.
    \item \texttt{Author}: l'autore del plugin.
    \item \texttt{Author URI}: URL dell'autore.
\end{itemize}

\section{Percorsi e URL}

Durante lo sviluppo di un plugin si ha spesso bisogno di riferirsi ad altri file all'interno della struttura del plugin stesso.

\subsection{Percorsi (paths)}

I percorsi vengono usati per \textbf{includere file PHP} dalla struttura del plugin. La funzione \texttt{plugin\_dir\_path()} restituisce il percorso sul file system:

\begin{lstlisting}[language=PHP]
define('MY_PLUGIN_PATH', plugin_dir_path(__FILE__));
// risultato: /home/nomeutente/sito.com/public_html/
//            wp-content/plugins/my_plugin/
\end{lstlisting}

Per includere un file si usa:

\begin{lstlisting}[language=PHP]
require_once(MY_PLUGIN_PATH . 'framework/MyClass.php');
\end{lstlisting}

\subsection{URL}

Gli URL vengono usati per referenziare risorse che devono essere caricate dal browser: \textbf{CSS, JavaScript, immagini}. La funzione \texttt{plugin\_dir\_url()} restituisce l'URL:

\begin{lstlisting}[language=PHP]
define('MY_PLUGIN_URL', plugin_dir_url(__FILE__));
// risultato: http://sito.com/wp-content/plugins/my_plugin/
\end{lstlisting}

Esempio di registrazione di uno script JavaScript:

\begin{lstlisting}[language=PHP]
wp_register_script(
    'my',
    MY_PLUGIN_URL . 'js/my.js',
    array('jquery'),
    '1.0',
    true
);
\end{lstlisting}

\section{Azioni e filtri (Hook)}

WordPress utilizza un sistema di \emph{hook} (``agganci'') per permettere ai plugin di interagire con il funzionamento del CMS:

\begin{itemize}
    \item Le \textbf{action} (azioni) vengono invocate durante un evento di WordPress. Un plugin può registrare una funzione da eseguire quando si verifica un certo evento.
    \item I \textbf{filter} (filtri) servono a modificare l'output restituito dai componenti dell'applicazione. Un plugin può intercettare e modificare un dato prima che venga utilizzato da WordPress.
\end{itemize}

Le funzioni principali per registrare gli hook sono:

\begin{itemize}
    \item \texttt{add\_action()}: registra una funzione da eseguire quando si verifica una certa action.
    \item \texttt{add\_filter()}: registra una funzione per modificare un certo dato tramite un filter.
\end{itemize}

\section{Esempio: plugin ``Hello World''}

Costruiamo, passo dopo passo, un semplice plugin che stampa ``hello world'' nel footer delle pagine e offre un'interfaccia di amministrazione per modificare il testo.

\subsection{Header del plugin}

\begin{lstlisting}[language=PHP]
<?php
/**
 * Plugin Name: Test plugin
 * Plugin URI:  https://mainwp.com
 * Description: This plugin does some stuff with WordPress
 * Version:     1.0.0
 * Author:      Your Name Here
 * Author URI:  https://mainwp.com
 * License:     GPL2
 */
?>
\end{lstlisting}

\subsection{Uso di add\_action()}

Aggiungiamo una funzione che stampa ``hello world'' nel footer della pagina. L'action \texttt{wp\_footer} viene invocata da WordPress al momento del rendering del footer di ogni pagina:

\begin{lstlisting}[language=PHP]
add_action('wp_footer', 'my_function');

function my_function() {
    echo 'hello world';
}
\end{lstlisting}

In questo modo, ogni volta che WordPress esegue l'hook \texttt{wp\_footer}, verrà chiamata la funzione \texttt{my\_function()} che stampa il testo.

\subsection{Creare una voce nel menu di amministrazione}

Per permettere all'amministratore di configurare il plugin dalla bacheca, creiamo una voce nel menu di amministrazione. La funzione \texttt{add\_management\_page()} aggiunge una pagina sotto il menu ``Strumenti'':

\begin{lstlisting}[language=PHP]
function my_admin_menu() {
    add_management_page(
        'Footer Text',         // titolo della pagina
        'Footer Text',         // titolo del menu
        'manage_options',      // capability richiesta
        __FILE__,              // slug
        'footer_text_admin_page' // funzione di callback
    );
}
\end{lstlisting}

\subsection{Modificare l'interfaccia di amministrazione}

Ecco la funzione di callback che costruisce l'interfaccia della pagina di amministrazione del plugin. Quando l'utente invia il form, il valore viene salvato nella tabella \texttt{wp\_options} tramite \texttt{update\_option()}. Per recuperare il valore si usa \texttt{get\_option()}:

\begin{lstlisting}[language=PHP]
function footer_text_admin_page() {
    $textvar = get_option('test_plugin_variable', 'hello world');

    if (isset($_POST['change-clicked'])) {
        update_option('test_plugin_variable',
                      $_POST['footertext']);
        $textvar = get_option('test_plugin_variable',
                              'hello world');
    }
    ?>
    <div class="wrap">
        <h1>Footer Text</h1>
        <p>This simple plugin will output some text to the
        footer of your template. Change the text below:</p>
        <form action="<?php echo $_SERVER['REQUEST_URI']; ?>"
              method="post">
            Footer Text:
            <input type="text"
                   value="<?php echo $textvar; ?>"
                   name="footertext"
                   placeholder="hello world"><br/>
            <input name="change-clicked"
                   type="hidden" value="1" />
            <input type="submit" value="Change Text" />
        </form>
    </div>
    <?php
}
\end{lstlisting}

\subsection{Commento sul codice}

Il codice mostrato illustra i meccanismi fondamentali di un plugin WordPress:

\begin{itemize}
    \item Uso degli \textbf{hook} (\texttt{add\_action}) per collegare una funzione a un evento di WordPress.
    \item Uso di funzioni come \texttt{get\_option()} e \texttt{update\_option()} per memorizzare e recuperare dati persistenti nel database di WordPress.
    \item Creazione di una \textbf{pagina di amministrazione} nel backend tramite \texttt{add\_management\_page()}.
    \item Utilizzo di un \textbf{form HTML} per permettere all'amministratore di modificare il testo.
\end{itemize}

\section{Considerazioni finali}

Lo sviluppo di plugin WordPress segue pattern ben definiti, basati sul sistema di hook (action e filter). La conoscenza di PHP, HTML e dei principi di programmazione web è un prerequisito; inoltre è importante consultare la documentazione ufficiale di WordPress per le API disponibili:

\begin{itemize}
    \item \url{https://developer.wordpress.org/plugins/}
    \item \url{https://developer.wordpress.org/rest-api/}
    \item \url{https://developer.wordpress.org/reference/}
\end{itemize}

Seguire le best practice per la sicurezza (validazione dell'input, sanitizzazione dell'output, uso di nonce per la protezione CSRF) è fondamentale, soprattutto nello sviluppo di plugin destinati alla distribuzione pubblica.
